% Part: first-order-logic
% Chapter: axiomatic-deduction
% Section: deduction-theorem

% verification of properties of provability needed for maximally
% consistent sets in the completeness chapter.

\documentclass[../../../include/open-logic-section]{subfiles}

\begin{document}

\iftag{FOL}
      {\olfileid{fol}{axd}{ded}}
      {\olfileid{pl}{axd}{ded}}

\olsection{નિગમન પ્રમેય}

આપણે જોયું તેમ, સ્વયંસિદ્ધિમૂલક તંત્રમાં !!{derivation}s આપવી કંટાળાજનક
છે અને !!{derivation}s શોધવી પણ મુશ્કેલ હોઈ શકે છે. !!{formula}sની લાંબી
યાદીઓ ખરેખર લખવાને બદલે, થોડાં સરળ પરિણામો વાપરીને એવી
!!{derivation}sનું અસ્તિત્વ સ્થાપવું સામાન્ય રીતે સહેલું છે. એવાં ત્રણ
પરિણામો આપણે પહેલેથી સ્થાપ્યાં છે: $!A \in \Gamma$ જાણતા હોઈએ ત્યારે
$\Gamma \Proves !A$ કહી શકાય એવું \olref[ptn]{prop:reflexivity} કહે છે.
જો $\Gamma \Proves !A$, તો $\Gamma \cup \{!B\} \Proves !A$ પણ થાય
એવું \olref[ptn]{prop:monotonicity} કહે છે. અને જો
$\Gamma \Proves !A$ તથા $!A \Proves !B$, તો $\Gamma \Proves !B$ એવું
\olref[ptn]{prop:transitivity}માંથી ફલિત થાય છે. હવે મોડસ પોનેન્સનું
``અધિ''-સ્વરૂપ આપતું બીજું સરળ પરિણામ જોઈએ:

\begin{prop}
  \ollabel{prop:mp} જો $\Gamma \Proves !A$ અને $\Gamma \Proves !A \lif
  !B$, તો $\Gamma \Proves !B$.
\end{prop}

\begin{proof}
  આપણી પાસે $\{!A, !A \lif !B\} \Proves !B$ છે:
  \begin{derivation}
    1. & $!A$ & પૂર્વ.\\
    2. & $!A \lif !B$ & પૂર્વ.\\
    3. & $!B$ & 1, 2, \MP
  \end{derivation}
  \olref[ptn]{prop:transitivity} વડે $\Gamma \Proves !B$.
\end{proof}

આ સંદર્ભમાં આપણે વાપરીશું તે સૌથી મહત્વનું પરિણામ નિગમન પ્રમેય છે:

\begin{thm}[નિગમન પ્રમેય]
\ollabel{thm:deduction-thm} $\Gamma \cup \{!A\} \Proves !B$ જો અને
માત્ર જો $\Gamma \Proves !A \lif !B$.
\end{thm}

\begin{proof}
``જો'' દિશા તરત મળે છે. જો $\Gamma \Proves !A \lif !B$, તો
\olref[ptn]{prop:monotonicity} વડે $\Gamma \cup \{!A\} \Proves !A \lif
!B$ પણ થાય. વળી, \olref[ptn]{prop:reflexivity} વડે
$\Gamma \cup \{!A\} \Proves !A$. તેથી \olref{prop:mp} વડે
$\Gamma \cup \{!A\} \Proves !B$.

``માત્ર જો'' દિશા માટે આપણે~$\Gamma \cup \{!A\}$માંથી $!B$ની
!!{derivation}ની લંબાઈ પર ગણિતીય અનુમાન કરીએ છીએ.

આધાર-કિસ્સામાં આપણે લંબાઈ~$1$ની દરેક !!{derivation} માટે દાવો સાબિત
કરીએ છીએ. $\Gamma \cup \{!A\}$માંથી $!B$ની લંબાઈ~$1$ની
!!^a{derivation} માત્ર~$!B$થી બને છે; અને તે યોગ્ય હોય તો $!B$ કાં તો
$\Gamma \cup \{!A\}$માં છે અથવા સ્વયંસિદ્ધિ છે. જો $!B \in \Gamma$
અથવા સ્વયંસિદ્ધિ હોય, તો $\Gamma \Proves !B$. વળી,
\olref[prp]{ax:lif1} વડે $\Gamma \Proves !B \lif (!A \lif !B)$ થાય,
અને \olref{prop:mp}થી $\Gamma \Proves !A \lif !B$ મળે. જો
$!B \in \{!A\}$, તો $\Gamma \Proves !A \lif !B$, કારણ કે છેલ્લું
!!{sentence}~$!A \lif !B$ એ $!A \lif !A$ જ છે અને તેને આપણે
\olref[pro]{ex:identity}માં !!{derive}d કર્યું છે.

\sourcecorrection{OLAX-017}{મૂળની \texttt{deduction-theorem.tex},
પંક્તિઓ 66--67માં સભ્યતા-ચિહ્નને તેના ડાબા પદાર્થ~$!B$થી જુદા
ગણિતીય ખંડમાં મૂકી “$!B$ is either $\in\Gamma\cup\{!A\}$” લખાયું
છે. અહીં સુનિર્ધારિત સભ્યતા-દાવો $!B\in\Gamma\cup\{!A\}$ વ્યક્ત કર્યો
છે.}

અનુમાનના પગલામાં ધારો કે~$\Gamma \cup \{!A\}$માંથી $!B$ની
!!a{derivation}નું છેલ્લું પગલું~$!B$ મોડસ પોનેન્સ વડે પ્રમાણિત થાય છે.
(જો તે મોડસ પોનેન્સ વડે પ્રમાણિત ન થતું હોય, તો $!B \in \Gamma$,
$!B \ident !A$, અથવા $!B$ સ્વયંસિદ્ધિ છે અને આધાર-કિસ્સાનું એ જ તર્ક
લાગુ પડે છે.) ત્યારે કોઈ !!{formula}~$!C$ માટે !!{derivation}નાં અગાઉનાં
પગલાં $!C \lif !B$ અને $!C$ છે; એટલે
$\Gamma \cup \{!A\} \Proves !C \lif !B$ અને
$\Gamma \cup \{!A\} \Proves !C$. તેમની !!{derivation}s ટૂંકી હોવાથી
તેમને અનુમાનની પૂર્વધારણા લાગુ પડે છે. તેથી આપણી પાસે બંને છે:
\begin{align*}
  & \Gamma \Proves !A \lif (!C \lif !B); \\
  & \Gamma \Proves !A \lif !C.
\end{align*}
પરંતુ
\[
\Gamma \Proves (!A \lif (!C \lif !B)) \lif
((!A\lif !C)  \lif (!A \lif !B)),
\]
એ પણ \olref[prp]{ax:lif2} વડે મળે છે; હવે \olref{prop:mp}ના બે ઉપયોગ
જરૂર મુજબ $\Gamma \Proves !A \lif !B$ આપે છે.
\end{proof}

ધ્યાન આપો કે નિગમન પ્રમેય સાચું રહે તે માટે જ
\olref[prp]{ax:lif1} અને \olref[prp]{ax:lif2}ની પસંદગી ચોક્કસ આ રીતે
કરાઈ હતી.

નીચેનાં !!{derivability} વિશેનાં કેટલાંક ઉપયોગી તથ્યો છે; તેમની
સાબિતી પ્રશ્ન તરીકે રાખીએ છીએ.

\begin{prop}
  \ollabel{prop:derivfacts}
\begin{enumerate}
\item $\Proves (!A \lif !B) \lif ((!B \lif !C)
  \lif (!A \lif !C))$; \ollabel{derivfacts:a}
\item જો $\Gamma \cup \{ \lnot !A\}
  \Proves \lnot !B$, તો $\Gamma \cup \{ !B\} \Proves
  !A$ (પ્રતિપક્ષન); \ollabel{derivfacts:b}
\item  $\{ !A, \lnot!A\} \Proves
    !B$ (એક્સ ફાલ્સો ક્વોડલિબેટ, વિસ્ફોટ); \ollabel{derivfacts:c}
\item  $\{ \lnot\lnot!A\} \Proves
  !A$ (દ્વિ-નિષેધ નિકાલ);\ollabel{derivfacts:d}
\item જો $\Gamma \Proves \lnot\lnot!A$, તો $\Gamma \Proves
  !A$;\ollabel{derivfacts:e}
\end{enumerate}
\end{prop}

\sourcecorrection{OLAX-005}{મૂળની \texttt{deduction-theorem.tex},
પંક્તિઓ 106--107માં
$\Proves (!A \lif !B) \lif ((!B \lif !C) \lif (!A \lif !C)$નું
બાહ્ય જમણું કૌંસ ખૂટે છે. અહીં સંતુલિત સૂત્ર મેળવવા તે કૌંસ ઉમેર્યું છે.}

\tagprob{FOL}
\begin{prob}
  \olref[fol][axd][ded]{prop:derivfacts} સાબિત કરો.
\end{prob}
\tagendprob

\tagprob{notFOL}
\begin{prob}
  \olref[pl][axd][ded]{prop:derivfacts} સાબિત કરો.
\end{prob}
\tagendprob

\end{document}
